\documentclass[a4paper,landscape]{article}

% Compile with LuaLaTeX or XeLaTeX (fontspec does not work with pdfLaTeX).
% With latexmk use:  latexmk -lualatex -pvc main.tex

\usepackage[margin=0.09cm, top=0.10cm, bottom=0.10cm]{geometry}
\usepackage{fontspec}
\usepackage{multicol}
\usepackage{multicolrule}
\usepackage{microtype}
\usepackage{xcolor}
\usepackage{enumitem}
\usepackage{amsmath}
\usepackage{amssymb}
\usepackage{listings}
\usepackage{graphicx}
\usepackage{tcolorbox}
\tcbuselibrary{raster,listings,skins,breakable}

\setlength{\columnseprule}{0.4pt}
\setlength{\columnsep}{0.16cm}
\pagestyle{empty}
\setlength{\parindent}{0pt}
\setlength{\parskip}{0pt}
\setlength{\multicolsep}{0pt}
\setlist{nosep}

% --- Configurable font variables ---
\newcommand{\cheatsheetfont}{Charter Bold}
\newcommand{\codefontname}{IBM Plex Mono}

\IfFontExistsTF{\cheatsheetfont}{%
  \newfontfamily\bodyfont{\cheatsheetfont}%
}{%
  % Portable fallback for systems without Charter Bold.
  \newfontfamily\bodyfont{Noto Serif Bold}%
}

\IfFontExistsTF{\codefontname}{%
  \newfontfamily\codefont{\codefontname}[Scale=0.85, Ligatures=TeXOff, Renderer=Basic]%
}{%
  % Portable fallback for systems without IBM Plex Mono.
  \newfontfamily\codefont{Noto Sans Mono}[Scale=0.85, Ligatures=TeXOff, Renderer=Basic]%
}

% Islands Light (RustRover default light) syntax colors
\definecolor{rustText}{HTML}{080808}
\definecolor{rustKeyword}{HTML}{0033B3}
\definecolor{rustPrimitiveType}{HTML}{0033B3}
\definecolor{rustStdType}{HTML}{336ECC}
\definecolor{rustUserType}{HTML}{080808}
\definecolor{rustFunction}{HTML}{00627A}
\definecolor{rustString}{HTML}{067D17}
\definecolor{rustNumber}{HTML}{1750EB}
\definecolor{rustComment}{HTML}{006400}
\definecolor{rustMacro}{HTML}{7A3DB8}
\definecolor{rustPunct}{HTML}{6C707E}
\definecolor{codebg}{HTML}{FFFFFF}

% Rust listings setup
\lstdefinelanguage{RustIsland}{%
  morekeywords={%
    as, async, await, break, const, continue, crate, dyn, else, enum, extern,
    false, fn, for, if, impl, in, let, loop, match, mod, move, mut, pub, ref,
    return, self, Self, static, struct, super, trait, true, type, unsafe, use,
    where, while%
  },
  sensitive=true,
  morecomment=[l]{//},
  morecomment=[s]{/*}{*/},
  morestring=[b]",
  morestring=[b]',
  moredelim=[is][\color{rustPrimitiveType}]{'}{'},
}[keywords,comments,strings]

\lstset{
  language=RustIsland,
  basicstyle=\codefont\color{rustText}\fontsize{2.88}{2.98}\selectfont,
  keywordstyle=\color{rustKeyword},
  commentstyle=\color{rustComment}\itshape,
  stringstyle=\color{rustString},
  numberstyle=\color{rustNumber},
  identifierstyle=\color{rustText},
  backgroundcolor=\color{codebg},
  showstringspaces=false,  
  breaklines=false,
  breakatwhitespace=false,
  tabsize=2,
  columns=fullflexible,
  keepspaces=true,
  lineskip=-0.40pt,
  xleftmargin=0pt,
  framexleftmargin=0pt,  
  aboveskip=0pt,
  belowskip=0pt,
  escapeinside={(*@}{@*)},
  literate={*}{{\raisebox{0.25ex}{\texttt{*}}}}1,
}

% Token classes (listings can't auto-detect everything, so we tag common names)
\lstset{
  % functions / methods
  emph={top_k_frequent, new, entry, or_default, values, max, copied, rev, get, extend, iter, push, len, cloned, as_ref, map, unwrap_or, range, sum, collect, println, print, eprintln, format, panic, assert, dbg, Self, is_none, is_some, unwrap, min, sort_by, cmp, swap, fill, reverse, push_back, pop_front, pop},
  emphstyle=\color{rustFunction},
  % macros / attributes
  emph={[2]vec, derive, test, cfg, allow, warn, deny, macro_rules, inline, must_use},
  emphstyle={[2]\color{rustMacro}},
  % primitive types
  emph={[3]i8, i16, i32, i64, i128, isize, u8, u16, u32, u64, u128, usize, f32, f64, bool, char, str},
  emphstyle={[3]\color{rustPrimitiveType}},
  % std library types and enum variants
  emph={[4]Vec, HashMap, HashSet, BTreeMap, BTreeSet, Option, Result, Box, String, Some, None, Ok, Err, BinaryHeap, VecDeque, Reverse, Graph, WGraph, EdgeList, MinHeap},
  emphstyle={[4]\color{rustStdType}},
  % user-defined types (add your own here)
  emph={[5]Solution, UnionFind, TreeNode, DSU, Edge},
  emphstyle={[5]\color{rustUserType}},
}

% Charter Bold is already the document font, so \textbf adds no visual distinction.
% Redefine it to underline instead, which remains visible when the whole document is bold.
\renewcommand{\textbf}[1]{\underline{#1}}

% Type scale: body < subsection < section, all in Charter Bold
\newcommand{\cheatsize}{\fontsize{3.28}{3.38}\selectfont}
\newcommand{\subsectionsize}{\fontsize{3.92}{4.02}\selectfont}
\newcommand{\sectionsize}{\fontsize{4.55}{4.65}\selectfont}

% Title styling
% Document title: centered, bold, no rule
\newcommand{\doctitle}[1]{%
  \par\vspace{0.03cm}%
  {\centering\sectionsize\bfseries #1\par}%
  \vspace{0.04cm}%
}
\newcounter{cheatsection}
% Major section: thin horizontal rule above, bold title, tiny space below
\newcommand{\sectiontitle}[1]{%
  \par\vspace{0pt}%
  \noindent{\rule{\linewidth}{0.25pt}}\par%
  \vspace{0pt}%
  \refstepcounter{cheatsection}%
  {\sectionsize\bfseries \arabic{cheatsection}.\ #1}\par%
  \vspace{0pt}%
}

\newcommand{\subsectiontitle}[1]{%
  \par\vspace{0pt}%
  {\subsectionsize\bfseries\itshape\textbf{#1}}\par%
  \vspace{0pt}%
}
% Proof bank: compact labeled proof cards, less wall-of-text than paragraphs
\newcommand{\proofbanktitle}[1]{%
  \par\vspace{0pt}%
  {\subsectionsize\bfseries\itshape\color{rustKeyword} #1}\par%
  \vspace{0pt}%
}
\newcommand{\proofitem}[2]{%
  \par\vspace{0pt}%
  \noindent\hangindent=0.42cm\hangafter=1%
  {\bfseries\color{rustFunction} #1.}\hspace{0.055cm}#2\par%
}

% Flex-like local wrapping boxes inside one main cheatsheet column.
% Usage:
% \begin{flexboxes}{3}
%   \flexbox{...}
%   \flexbox{...}
%   \flexbox{...}
% \end{flexboxes}
\newenvironment{flexboxes}[1]{%
  \begin{tcbraster}[
    raster width=\linewidth,
    raster columns=#1,
    raster equal height=none,
    raster column skip=0.010cm,
    raster row skip=0.006cm,
    raster left skip=0pt,
    raster right skip=0pt
  ]
}{%
  \end{tcbraster}
}

\newcommand{\flexbox}[1]{%
  \begin{tcolorbox}[
    enhanced,
    sharp corners,
    boxrule=0pt,
    colback=white,
    left=0pt,right=0pt,top=0pt,bottom=0pt,
    boxsep=0pt,
    width=\linewidth,
    before skip=0pt,after skip=0pt
  ]#1\end{tcolorbox}%
}

\newcommand{\flexcode}[1]{%
  \begin{tcolorbox}[
    enhanced,
    sharp corners,
    boxrule=0pt,
    colback=white,
    left=0pt,right=0pt,top=0pt,bottom=0pt,
    boxsep=0pt,
    width=\linewidth,
    before skip=0pt,after skip=0pt
  ]\begin{lstlisting}#1\end{lstlisting}\end{tcolorbox}%
}
% Local code columns with real vertical separators. Works better than minipage+\vrule.
\newcommand{\codevsep}{\vrule width 0.25pt}

\newcommand{\twocodecols}[2]{%
\noindent\begin{tabular}{@{}p{0.485\linewidth}@{\hspace{0.01\linewidth}\codevsep\hspace{0.01\linewidth}}p{0.485\linewidth}@{}}%
#1 & #2
\end{tabular}%
}

\newcommand{\threecodecols}[3]{%
\noindent\begin{tabular}{@{}p{0.30\linewidth}@{\hspace{0.01\linewidth}\codevsep\hspace{0.01\linewidth}}p{0.30\linewidth}@{\hspace{0.01\linewidth}\codevsep\hspace{0.01\linewidth}}p{0.30\linewidth}@{}}%
#1 & #2 & #3
\end{tabular}%
}
\usepackage{array}

% Worked exam problem cards and compact tags.
\newcommand{\workedq}[3]{%
  \par\vspace{0pt}%
  \noindent{\color{rustKeyword}\bfseries #1. #2}\par\vspace{-0.006cm}%
  \noindent\hangindent=0.22cm\hangafter=1 #3\par%
}
\newcommand{\wk}[1]{{\color{rustMacro}\bfseries #1}}
\newcommand{\ans}{\wk{Ans:} }
\newcommand{\inv}{\wk{Inv.:} }
\newcommand{\init}{\wk{Init:} }
\newcommand{\maint}{\wk{Maint:} }
\newcommand{\term}{\wk{Term:} }


\begin{document}
\cheatsize
\bodyfont
\linespread{0.70}\selectfont
\setlength{\tabcolsep}{1.0pt}
\renewcommand{\arraystretch}{0.70}
\setlength{\fboxsep}{0pt}
\setlength{\fboxrule}{0pt}
\sloppy
\setlength{\abovedisplayskip}{0pt}
\setlength{\belowdisplayskip}{0pt}
\setlength{\abovedisplayshortskip}{0pt}
\setlength{\belowdisplayshortskip}{0pt}
\setlist[itemize]{nosep,topsep=0pt,parsep=0pt,partopsep=0pt,leftmargin=*,itemsep=0pt}

\begin{multicols*}{4}
\flushcolumns
\SetMCRule{line-style=solid, extend-fill}

\sectiontitle{Complexity Table / Fast Lookup}

\resizebox{\linewidth}{!}{%
\begin{tabular}{@{}llll@{}}
\hline
Algorithm/Data Struct & Time & Space & Notes \\
\hline
Hash search/insert/delete & expected $O(1)$, worst $O(n)$ & $O(n)$ & $\alpha$=n/m \\
BST ops & $O(h)$ & $O(n)$ & h $\log n$ balanced, n worst \\
Inorder BST & $\Theta(n)$ & $O(h)$ & sorted output \\
Insertion sort & best $\Theta(n)$, worst $\Theta(n^2)$ & $O(1)$ & stable \\
Merge sort & $\Theta(n\;\log\;n)$ & $\Theta(n)$ & stable \\
Quicksort randomized & expected $\Theta(n\;\log\;n)$ & $O(\log n)$ & worst $\Theta(n^2)$ \\
Counting sort & $\Theta(n+k)$ & $\Theta(n+k)$ & integer keys \\
Heapify & $O(\log n)$ & recursion & assumes children heaps \\
Build heap & $\Theta(n)$ & $O(1)$ & bottom-up \\
Heapsort & $\Theta(n\;\log\;n)$ & $O(1)$ & not stable \\
PQ insert/extract & $O(\log n)$ & $O(n)$ & heap \\
BFS & $\Theta(V+E)$ & $\Theta(V)$ & unweighted shortest paths \\
DFS & $\Theta(V+E)$ & $\Theta(V)$ & topo/cycles/SCC \\
Topological sort & $\Theta(V+E)$ & $\Theta(V)$ & DAG only \\
Kosaraju SCC & $\Theta(V+E)$ & $\Theta(V+E)$ & 2 DFS \\
Union-Find op & $O(\alpha(n))$ & $O(n)$ & rank+compression \\
Kruskal & $O(E\log E)$ & $O(V)$ & sort + DSU \\
Prim & $O(E\log V)$ & $O(V+E)$ & PQ \\
Dijkstra & $O(E\log V)$ & $O(V+E)$ & no negative edges \\
Bellman-Ford & $\Theta(VE)$ & $O(V)$ & negative edges OK \\
DAG shortest path & $\Theta(V+E)$ & $O(V)$ & topo order \\
Ford-Fulkerson & $O(E|f^*|)$ & $O(V^2{+}E)$ & integer caps, residual \\
Edmonds-Karp & $O(VE^2)$ & $O(V^2{+}E)$ & BFS residual \\
Rod cutting & $\Theta(n^2)$ & $\Theta(n)$ & DP \\
Matrix chain & $\Theta(n^3)$ & $\Theta(n^2)$ & DP \\
LCS & $\Theta(nm)$ & $\Theta(nm)$ & DP \\
Knapsack 0/1 & $\Theta(nW)$ & $\Theta(nW)/\allowbreak \Theta(W)$ & pseudo-poly \\
\hline
Coin change & $\Theta(W\cdot \#\text{coins})$ & $\Theta(W)$ & unbounded \\
\hline
\end{tabular}%
}
\proofbanktitle{Exam Pattern Matcher}
\begin{tabular}{@{}ll@{}}
\textbf{new optimization + choices} & DP: state/base/transition/order/answer/proof \\
\textbf{local best + exchange} & greedy: safe choice + induction \\
\textbf{connect all min weight} & MST: Kruskal/Prim + cut property \\
\textbf{unweighted shortest path} & BFS levels \\
\textbf{weighted no negative} & Dijkstra \\
\textbf{negative weights / cycle} & Bellman--Ford \\
\textbf{dependencies / precedence} & DAG + topological order \\
\textbf{capacity / matching / disjoint paths} & max-flow reduction \\
\textbf{expected count} & indicators + linearity \\
\textbf{online unknown future} & competitive ratio vs OPT \\
\textbf{min max feasible value} & binary search answer + monotone feasible \\
\textbf{mutual reachability forbidden} & SCC condensation / comp[u]=comp[v] \\
\textbf{capacitated assignment} & max-flow with integral capacities \\
\textbf{tree path/subtree optimum} & postorder DP: inside child / through root \\
\textbf{tree traversals reconstruct} & inorder+post/pre split at root \\
\textbf{time-dependent edges} & earliest-arrival Dijkstra/events \\
\textbf{grid/assignment/paths} & bipartite matching / max-flow
\end{tabular}

\proofbanktitle{Exam Writeup Checklist}
\textbf{Algorithm:} input/output, data structs, return value.
\textbf{Correctness:} complete invariant/induction/exchange proof, not idea.
\textbf{DP:} State, Base, Transition, Order, Answer, Runtime, Space.
\textbf{Bounds:} define $n,V,E,W$; justify loops.
\textbf{Conventions:} $<$ vs $\le$, inclusive/exclusive, 0/1-index, reachable vs global negative cycle.

\proofbanktitle{Proof Bank: Generic Templates}
\proofitem{Loop inv}{Init true before loop. Maint: one step preserves. Term: invariant + stop condition $\Rightarrow$ answer.}

\proofitem{Induction/substitution}{Claim upper/lower bound with constants. State IH for all smaller inputs. Substitute IH. Collect terms. Pick constants so inequality holds. Base cases: choose $c$ large enough for all $n<n_0$.}

\proofitem{Substitution trap}{If $T(n)\le c\sum_i(a_in)^p+bn$ and $\sum_i a_i^p=1$, then $cn^p+bn$ FAILS. No slack absorbs $+bn$. Strengthen: $T(n)\le cn^p-dn^q$ with lower-order subtraction.}

\proofitem{Midterm recurrence}{$T(n)=T(4n/5)+T(3n/5)+an$. Guess $\Theta(n^2)$ because $(4/5)^2+(3/5)^2=1$. Lower: $T(n)\ge T(4n/5)+T(3n/5)\ge c((4n/5)^2+(3n/5)^2)=cn^2$. Upper must be strengthened: assume $T(m)\le cm^2-dm$ for all $m<n$. Then $T(n)\le c(4n/5)^2-d(4n/5)+c(3n/5)^2-d(3n/5)+an=cn^2-7dn/5+an\le cn^2-dn$ iff $d\ge5a/2$. Choose $d$ so, then choose $c$ large for base cases.}

\proofitem{D\&C}{Induct on size. Recursive calls solve smaller problems. Combine covers all optimum cases and constructs/chooses correct result.}

\proofitem{DP}{State meaning. Base correct. Induct in fill order. Transition enumerates all valid first/last choices; dependencies already optimal. Optimal solution uses one enumerated choice; every candidate feasible. Therefore $dp[s]$ optimum. Write answer $=dp[\dots]$.}

\proofitem{Parametric search}{Guess answer $M$. Build feasible$(M)$. Prove monotone: if $M$ works then any larger $M$ works (minimization). Binary search smallest feasible $M$. Correctness = binary search invariant + feasible predicate exact.}

\proofitem{Greedy}{Safe choice: for optimum $O$ containing current partial solution, if chosen $e\notin O$, exchange $f\in O$ with $e$ without worsening. Then induct.}

\proofbanktitle{Proof Bank: Sorting / Hash / DS}
\proofitem{BST inorder}{Induct on subtree. Left keys $<$ root, right keys $\ge$ root. IH sorts left/right. Left,root,right sorted and complete.}

\proofitem{Rank/interval}{Store $size=1+size(L)+size(R)$. Rank follows one search path; when going right add $size(L)+1$. Thus $O(h)$. $[a,b]=rank_{\le}(b)-rank_{<}(a)$.}

\proofitem{Heapify}{Children are heaps. Swap with larger child if violated; parent fixed. Only swapped child may violate; recurse to leaf. Cost height $O(\log n)$.}

\proofitem{BuildHeap}{Leaves are heaps. Bottom-up: children already heaps, so Heapify fixes subtree. Runtime: nodes height $h\le n/2^{h+1}$, total $\sum_h O(h)n/2^h=O(n)$.}

\proofitem{Sorting lower bound}{Comparison sort decision tree needs $\ge n!$ leaves. Worst-case comparisons $\ge$ height $\ge\log_2(n!)=\Omega(n\log n)$.}

\proofitem{Quicksort exp}{Pair compared at most once. $X_{ij}=1$ if ranks $i,j$ compared. Compare iff first pivot among ranks $i..j$ is $i$ or $j$, so $P(X_{ij})=2/(j-i+1)$. Sum $O(n\log n)$.}

\proofitem{Hashing}{Single key: for key $x$, $X_y=1[h(y)=h(x)]$. Universal hashing gives $E[X_y]\le1/m$, so expected chain collisions $\le(n-1)/m=\Theta(\alpha)$ and search $O(1+\alpha)$. Total collisions: $X=\sum_{i<j}X_{ij}$, $X_{ij}=1[h(x_i)=h(x_j)]$, hence $E[X]\le{n\choose2}/m=n(n-1)/(2m)=\Theta(n2/m)$. $k$-collision: $E[\#k\text{-collisions}]\le {n\choose k}/m^{k-1}$. Markov: $P[X\ge1]\le E[X]$.}
\proofitem{Hash mini}{Length-$\ell$ keys: pair collision prob $1/m^\ell$, so $E[\#collisions]={n\choose2}/m^\ell$; if $m^\ell\ge{n\choose2}$ then $P[\exists collision]\le1$. Random $q$-color vertices: edge mono prob $1/q$, so $E[\#mono]=|E|/q$.}
\proofitem{DSU}{Union by rank attaches smaller rank under larger. Rank increases only on equal-rank merge. Path compression preserves reps and shortens future paths. With both: $m$ ops in $O(m\alpha(n))$.}

\proofbanktitle{Proof Bank: Graph Traversal}
\proofitem{BFS}{Queue processed in nondecreasing distance. Discover $v$ from $u$: $d[v]=d[u]+1$. Shorter path would come from smaller-distance vertex processed earlier, so $v$ would already be discovered.}

\proofitem{DFS/topo}{Directed cycle iff DFS has back edge to GRAY ancestor. DAG has no back edges. For every edge $u\to v$, $f[u]>f[v]$; decreasing finish time is topo order.}

\proofitem{DFS edge facts}{Parenthesis theorem: finish intervals disjoint or nested; nested iff descendant. White-path theorem: $v$ descendant of $u$ iff when $u$ discovered, there is a white path $u\leadsto v$.}

\proofitem{SCC/Kosaraju}{Condensation graph is DAG. First DFS finish order puts source SCCs late. In reversed graph, DFS from highest finish cannot leave its SCC to an unvisited SCC. Each second-pass DFS outputs one SCC.}

\proofitem{SCC reachability pairs}{$u\leadsto v$ and $v\leadsto u$ iff $comp[u]=comp[v]$. For forbidden pair $(u,v)$ reject iff same component. Correct because SCC is maximal mutual reachability. Runtime $O(V+E)$.}

\proofitem{DAG SP/LP}{Topo order processes every predecessor before vertex. When processing $u$, $dist[u]$ already final over all paths ending at $u$; relaxing outgoing edges propagates optimum once. Runtime $\Theta(V+E)$.}

\proofbanktitle{Proof Bank: MST / Shortest Paths / Flow}
\proofitem{Cut property}{Let $e$ be light edge crossing cut respecting $A$. Take MST $T\supseteq A$ not using $e$. Add $e$ creates cycle with another crossing edge $f$. Since $w(e)\le w(f)$, $T-f+e$ is MST containing $A\cup\{e\}$.}

\proofitem{Cycle property}{If MST $T$ contains unique heaviest edge $e$ on cycle, remove $e$; cycle has another edge $f$ reconnecting cut with $w(f)<w(e)$. Then $T-e+f$ cheaper, contradiction.}

\proofitem{Kruskal}{DSU components define cut respecting chosen edges. Next accepted edge is lightest crossing some component cut. Cut property says safe. Induct until $|V|-1$ edges.}

\proofitem{Prim}{$S=$ vertices in tree. Heap chooses lightest edge crossing $(S,V\setminus S)$. Cut property says safe. Add one vertex; repeat.}

\proofitem{Dijkstra}{$S=$ finalized. Pick outside $x$ with min tentative dist. Any alternative path first crosses from $S$ to outside at $(u,v)$. Since $w\ge0$, path length $\ge d[u]+w(u,v)\ge d[v]\ge d[x]$. So $d[x]$ final.}

\proofitem{Bellman-Ford}{After $i$ passes, all shortest paths using $\le i$ edges are correct. Induct by last edge relaxation. If no negative cycle, shortest simple path has $\le V-1$ edges. If pass $V$ improves, reachable negative cycle.}

\proofitem{Max-flow}{If no augmenting path exists, let S be vertices reachable from s
in residual graph. For every edge $u\in S$,$v\notin S$: $cf(u,v)=0$, therefore $f(u,v)=c(u,v)$.
For every edge $u\notin S$,$v\in S$:
$cf(v,u)=0$, therefore $f(u,v)=0$.

Hence net flow crossing cut equals cut capacity:
$|f| = c(S,T)$.

Since every flow $\le$ every cut,
flow and cut are both optimal.}

\proofitem{Flow cert}{Any flow value $\le$ any cut capacity. Flow value $F$ + cut capacity $F$ proves both optimal.}

\proofitem{Bipartite matching flow}{Caps 1 force each left/right vertex used at most once. Integral capacities give integral max flow. Unit flow through $l\to r$ is matched pair; every matching gives same-size flow.}

\proofitem{Cap assignment flow}{Build $s\to$ people cap 1, person$\to$ acceptable item cap 1, item$\to t$ cap capacity. Feasible iff maxflow = \#people. Integral flow = assignment; assignment = integral flow.}

\proofitem{SP transform}{Multiplying every edge weight by $a>0$ preserves all shortest paths. Adding $a$ to every edge does not: paths with different number of edges get different extra cost.}

\proofitem{Resource SP}{State graph $(v,b)$ where $b=$ used budget / remaining battery. Add transition only if resource feasible. Run BFS/Dijkstra depending on edge costs. Size $O(VB+EB)$. Correct since every feasible walk corresponds to one state path and vice versa.}
\proofitem{2D sorted matrix}{Rows/cols sorted. Start top-right. If $A[i,j]>x$, move left; if $A[i,j]<x$, move down. Each step deletes one impossible row/col, so $O(n+m)$, $O(1)$.}
\proofbanktitle{Proof Bank: Trees / Exam Reductions}
\proofitem{Tree DP}{Postorder. Return height/best info from children. Best at $u$ is max(best in left, best in right, best crossing/using $u$). Induct on subtree: any optimal object either lies fully in one child or uses $u$.}

\proofitem{Tree diameter}{Let $h(u)=1+\max(h_L,h_R)$, empty height $-1$. $diam(u)=\max(diam_L,diam_R,h_L+h_R+2)$. Path either inside one child or crosses through $u$.}
\proofitem{Root path}{Longest simple path forced to pass through root: $h(root.left)+h(root.right)+2$ with empty height $-1$. Compute heights postorder in $O(n)$.}
\proofitem{Tree rebuild}{Inorder+postorder: root=last(post); split inorder at root; subtree sizes split postorder; recurse. Preorder+inorder: root=first(pre). Pre+post alone not unique.}

\proofitem{Unicyclic spanning trees}{Connected graph with $|E|=|V|$ has exactly one cycle $C$. Every spanning tree = delete exactly one edge of $C$. Deleting non-cycle edge disconnects; deleting cycle edge breaks only cycle. Hence \#spanning trees $=|C|$.}


\proofbanktitle{Proof Bank: Online / Random}

\proofitem{LRU paging}{Cache size $k$. Phase has $\le k$ distinct pages; new phase begins at request of $(k+1)$-st distinct. LRU: $\le k$ misses/phase, because requested page survives until $k$ other distinct pages requested. OPT: $\ge1$ miss/completed phase, because phase + first request of next phase contains $k+1$ distinct pages but cache holds $k$. Thus LRU $\le k\cdot$OPT$+O(k)$; LRU is $k$-competitive.}
\sectiontitle{Complexity / Asymptotics}

$\log n \ll \sqrt n \ll n \ll n\log n \ll n^2 \ll n^3 \ll 2^n \ll n! \ll n^n$

$O(g):f\le cg$ eventually. $\Omega(g):f\ge cg$. $\Theta(g)=O\cap\Omega$.
$o(g):f/g\to0$. $\omega(g):f/g\to\infty$.
Stirling: $n!\sim\sqrt{2\pi n}(n/e)^n$, so $\log(n!)=\Theta(n\log n)$.

\textbf{Useful sums:}

$\sum_{i=1}^n i=\frac{n(n+1)}2=\Theta(n^2),\;
\sum_{i=0}^{\log n}n=n\log n,\;
\sum_{i\ge0}2^{-i}=2,\;
\sum_{h\ge0}h2^{-h}=2$

\textbf{Master / recurrence solver:}

\textbf{Master:} $T(n)=aT(n/b)+f(n)$. Let $d=\log_b(a)$.
Compare $f(n)$ with $n^d$.

\begin{tabular}{@{}ll@{}}
$f$ smaller: $f=O(n^{d-\varepsilon})$ & $T=\Theta(n^d)$\\
same: $f=\Theta(n^d \log^k n)$ & $T=\Theta(n^d \log^{k+1}n)$\\
$f$ bigger: $f=\Omega(n^{d+\varepsilon})$ & $T=\Theta(f)$ if regularity\\
\end{tabular}

\textbf{Fast mental version if $f=n^k$:}
\begin{tabular}{@{}lll@{}}
$k<d$ & leaves win & $\Theta(n^d)$\\
$k=d$ & tie/layers equal & $\Theta(n^d \log n)$\\
$k>d$ & root/work wins & $\Theta(n^k)$\\
\end{tabular}

\textbf{Non-equal splits:} $T=\sum_i a_i T(b_i n)+\Theta(n^k)$.
Find $p$ s.t. $\sum_i a_i b_i^p = 1$.
\begin{tabular}{@{}lll@{}}
$p>k$ & leaves win & $\Theta(n^p)$\\
$p=k$ & tie & $\Theta(n^p \log n)$\\
$p<k$ & root/work wins & $\Theta(n^k)$\\
\end{tabular}

\textbf{Trap:} if $\sum_i a_i b_i^p=1$, plain $cn^p$ proof has no slack;
use $cn^p - d n^q$.
\textbf{Akra--Bazzi mini:} $T(n)=\alpha T(an)+\beta T(bn)+\Theta(n^d)$. Find $p$: $\alpha a^p+\beta b^p=1$. Then $p>d\Rightarrow\Theta(n^p)$, $p=d\Rightarrow\Theta(n^d\log n)$, $p<d\Rightarrow\Theta(n^d)$.
\textbf{Recurrences / runtimes:}
{\codefont\fontsize{3.15}{3.25}\selectfont
\begin{tabular}{@{}lll@{}}
binsearch & $T=T(n/2)+1$ & $\Theta(\log n)$\\
merge/maxsub & $T=2T(n/2)+n$ & $\Theta(n \log n)$\\
naive mm & $T=8T(n/2)+n^2$ & $\Theta(n^3)$\\
strassen & $T=7T(n/2)+n^2$ & $\Theta(n^{\log_2 7})$\\
heapify & $T\le T(2n/3)+1$ & $O(\log n)$\\
qsort worst & $T=T(n-1)+n$ & $\Theta(n^2)$\\
qsort bal & $T=2T(n/2)+n$ & $\Theta(n \log n)$\\
BF & $(V-1)E$ relax & $\Theta(V\cdot E)$\\
BFS/DFS & visit $V$, scan $E$ & $\Theta(V+E)$
\end{tabular}}

\textbf{Master map:}
\begin{tabular}{@{}ll@{}}
$T(n/2)+1 \to \log n$ & $T(n/2)+n \to n$\\
$2T(n/2)+1 \to n$ & $2T(n/2)+n \to n \log n$\\
$3T(n/2)+n \to n^{\log_2 3}$ & $7T(n/2)+n^2 \to$ Strassen\\
$8T(n/2)+n^2 \to n^3$ & $T(n-1)+n \to n^2$
\end{tabular}

\sectiontitle{Basic Data Structures}
\subsectiontitle{Stack / Queue / Linked List}
\begin{tabular}{@{}lll@{}}
Stack & push/pop/top & $O(1)$, array end, underflow if empty \\
Queue & enqueue/dequeue & $O(1)$ with circular array / deque \\
Linked list & search/insert/delete & search $O(n)$, known-node delete $O(1)$
\end{tabular}

Doubly linked list insert head: `x.next=head; if head!=NIL head.prev=x; head=x; x.prev=NIL`.
Delete known node: reconnect `prev.next` and `next.prev`; finding node still $O(n)$.

\subsectiontitle{Hash Tables}

\textbf{Idea:} bucket $h(key)\bmod m$. \textbf{Collisions:} chaining / probing. \textbf{Load factor} $\alpha=n/m$.

Hash ops: insert/delete expected O(1), search expected O(1+alpha), all worst O(n); m=Theta(n) => alpha=Theta(1).
\textbf{Expected total collisions:}
$X=\sum_{i<j}X_{ij}$, where $X_{ij}=1[h(x_i)=h(x_j)]$.
Universal hashing gives P[Xij=1]<=1/m, so E[X]=sum E[Xij] <= C(n,2)/m = n(n-1)/(2m)=Theta(n2/m).
For one key: $\mathbb E[\#\text{collisions}]\le(n-1)/m=\Theta(\alpha)$, hence expected search $O(1+\alpha)$.
\textbf{Traps:} HashMap/HashSet $O(1)$ is expected/amortized, not worst-case.
Bad hashing / all collisions $\Rightarrow O(n)$.

\textbf{k-collision stencil:} $X_S=1$ if fixed $k$-set collides. $P[X_S=1]\le1/m^{k-1}$. $E[X]=\sum_{|S|=k}E[X_S]\le {n\choose k}/m^{k-1}$. If need existence prob: $P[X\ge1]\le E[X]$.

\subsectiontitle{Binary Search Trees}
\textbf{Binary search (sorted array):}
\begin{lstlisting}
fn binary_search(a: Vec<_>,x){
  l=0; r=a.len();
  while l<r {
    m=(l+r)/2;
    if a[m]==x { return Some(m); }
    if a[m]<x  { l=m+1; }
    else       { r=m; }
  }
  None
}
//lower_bound = first >= x: First position where the value is >= x.
//upper_bound = first > x: First position where the value is > x.
\end{lstlisting}
\textbf{BST property:}
BST property: all keys in L(x) < x.key; all keys in R(x) >= x.key.

BST ops search/insert/delete/min/max/succ/pred O(h); inorder Theta(n); balanced h=Theta(log n), degenerate h=Theta(n).

% --- BST core operations + insert/transplant/delete: two minipages, no titles ---
\noindent
\begin{minipage}[t]{0.49\linewidth}\vspace{0pt}
\begin{lstlisting}[aboveskip=0pt,belowskip=0pt]
search(x,k):
  while x != NIL and x.key != k:
    if k < x.key: x = x.left
    else:         x = x.right
  return x

min(x): while x.left  != NIL: x = x.left;  return x
max(x): while x.right != NIL: x = x.right; return x

successor(x):
  if x.right != NIL: return min(x.right)
  y = x.parent
  while y != NIL and x == y.right:
    x = y
    y = y.parent
  return y
insert(T,z):
  y=NIL; x=T.root
  while x!=NIL:
    y=x
    if z.key<x.key: x=x.left else: x=x.right
  z.parent=y
  if y==NIL: T.root=z
  else if z.key<y.key: y.left=z
  else: y.right=z

\end{lstlisting}
\end{minipage}\hfill
\begin{minipage}[t]{0.49\linewidth}\vspace{0pt}
\begin{lstlisting}[aboveskip=0pt,belowskip=0pt]
transplant(T,u,v):
  if u.parent==NIL: T.root=v
  else if u==u.parent.left: u.parent.left=v
  else: u.parent.right=v
  if v!=NIL: v.parent=u.parent

delete(T,z):
  if z.left==NIL: transplant(T,z,z.right)
  else if z.right==NIL: transplant(T,z,z.left)
  else:
    y=min(z.right)
    if y.parent!=z:
      transplant(T,y,y.right); 
      y.right=z.right; y.right.parent=y
    transplant(T,z,y); 
    y.left=z.left; y.left.parent=y
\end{lstlisting}
\end{minipage}
\textbf{Delete cases:} no left, no right, or replace by successor. All $O(h)$.

\textbf{Traversals:}
\begin{lstlisting}[aboveskip=0pt,belowskip=0pt]
inorder(x):  if x!=NIL: inorder(x.left); visit(x); inorder(x.right)
preorder(x): if x!=NIL: visit(x); preorder(x.left); preorder(x.right)
postorder(x):if x!=NIL: postorder(x.left); postorder(x.right); visit(x)
\end{lstlisting}

inorder(BST)=sorted order.
 \textbf{Augmented BST / rank queries:}
Store \texttt{size(x)=1+size(left)+size(right)}. Rank follows one root-to-leaf path, so $O(h)$.
\texttt{rank\_<}(x)=\#keys $<x$, \texttt{rank\_<=}(x)=\#keys $\le x$.
Intervals: $[a,b]=\texttt{rank\_<=}(b)-\texttt{rank\_<}(a)$; $(a,b]=\texttt{rank\_<=}(b)-\texttt{rank\_<=}(a)$; $[a,b)=\texttt{rank\_<}(b)-\texttt{rank\_<}(a)$; $(a,b)=\texttt{rank\_<}(b)-\texttt{rank\_<=}(a)$.



\sectiontitle{Sorting Algorithms}

\subsectiontitle{Insertion Sort}
In-place, stable. Invariant: before iteration $j$, prefix $A[0..j)$ sorted.
\begin{lstlisting}
fn insertion_sort(a: &mut Vec<i32>) {
  for j in 1..a.len() {
    let key = a[j];
    let mut i = j;
    while i > 0 && a[i-1] > key {
      a[i] = a[i-1]; i -= 1;
    }
    a[i] = key;
  }
}
\end{lstlisting}
Best $\Theta(n)$, worst/avg $\Theta(n^2)$, space $O(1)$.
% --- Merge sort: two minipages, left = merge_sort, right = merge ---
\subsectiontitle{Merge Sort}
Stable, simple version uses $\Theta(n)$ temp.\\
\noindent
\begin{minipage}[t]{0.38\linewidth}\vspace{0pt}
\begin{lstlisting}[aboveskip=0pt,belowskip=0pt]
fn merge_sort(a,l,r){
  if r-l <= 1 { return; }
  let m=(l+r)/2;
  merge_sort(a,l,m);
  merge_sort(a,m,r);
  merge(a,l,m,r);
}
\end{lstlisting}
\end{minipage}\hfill
\begin{minipage}[t]{0.60\linewidth}\vspace{0pt}
\begin{lstlisting}[aboveskip=0pt,belowskip=0pt]
fn merge(a,l,m,r){
  let (mut i,mut j)=(l,m); let mut t=Vec::new();
  while i<m && j<r {
    if a[i] <= a[j] { t.push(a[i]); i+=1; }
    else { t.push(a[j]); j+=1; }
  }
  while i<m { t.push(a[i]); i+=1; }
  while j<r { t.push(a[j]); j+=1; }
  for k in 0..t.len(){ a[l+k]=t[k]; }
}
\end{lstlisting}
\end{minipage}
$T(n)=2T(n/2)+\Theta(n)=\Theta(n\log n)$.
% --- Quicksort: two minipages, left = partition, right = quicksort ---
\subsectiontitle{Quicksort}
In-place. Partition invariant: $A[l..i]\le p$, $A[i+1..j]>p$, unknown after $j$, pivot at $r$.
\noindent
\begin{minipage}[t]{0.49\linewidth}\vspace{0pt}
\begin{lstlisting}[aboveskip=0pt,belowskip=0pt]
fn partition(a,l,r)->usize{
  let p=a[r]; let mut i=l;
  for j in l..r {
    if a[j] <= p { a.swap(i,j); i+=1; }
  }
  a.swap(i,r); i
}
\end{lstlisting}
\end{minipage}\hfill
\begin{minipage}[t]{0.49\linewidth}\vspace{0pt}
\begin{lstlisting}[aboveskip=0pt,belowskip=0pt]
fn quicksort(a,l,r){
  if l>=r { return; }
  let q=partition(a,l,r);
  if q>0 { quicksort(a,l,q-1); }
  quicksort(a,q+1,r);
}
\end{lstlisting}
\end{minipage}
Balanced/expected randomized $\Theta(n\log n)$, worst $\Theta(n^2)$, stack expected $O(\log n)$ worst $O(n)$.

\subsectiontitle{Counting Sort / Linear Sorting}
Integer keys $0..k$, stable with prefix sums, not comparison-based.
\begin{lstlisting}
fn counting_sort(a:&Vec<usize>, k:usize)->Vec<usize>{
  let mut c=vec![0;k+1]; let mut out=vec![0;a.len()];
  for &x in a { c[x]+=1; }
  for i in 1..=k { c[i]+=c[i-1]; }
  for idx in (0..a.len()).rev(){
    let x=a[idx]; c[x]-=1; out[c[x]]=x;
  }
  out
}
\end{lstlisting}
$\Theta(n+k)$ time/space. Stable with prefix sums. Needs bounded integer keys.

\sectiontitle{Heaps / Priority Queues / Heapsort}

\subsectiontitle{Max-Heapify}
Assume children subtrees are heaps; only $A[i]$ may violate. 0-index:
\begin{lstlisting}
fn max_heapify(a:&mut Vec<i32>, i:usize, heap_size:usize){
  let (l,r)=(2*i,2*i+1); let mut largest=i;
  if l<heap_size && a[l]>a[largest] { largest=l; }
  if r<heap_size && a[r]>a[largest] { largest=r; }
  if largest!=i { a.swap(i,largest); max_heapify(a,largest,heap_size); }
}
\end{lstlisting}
$O(height(i))\le O(\log n)$. Trap: during heapsort, \texttt{heap\_size} $\ne$ \texttt{a.len()}.
\subsectiontitle{Build Max Heap}
\begin{lstlisting}
fn build_max_heap(a:&mut Vec<i32>){
  let n=a.len();
  for i in (0..(n/2)).rev(){ max_heapify(a,i,n); }
}
\end{lstlisting}
Tight: nodes height $h\le n/2^{h+1}$, cost $O(h)$, so $\sum_h nh/2^h=O(n)$; hence BuildHeap $\Theta(n)$, not $\Theta(n\log n)$.

\subsectiontitle{Heapsort}
\begin{lstlisting}
fn heapsort(a:&mut Vec<i32>){
  build_max_heap(a); let mut h=a.len();
  while h>1 { a.swap(0,h-1); h-=1; max_heapify(a,0,h); }
}
\end{lstlisting}
Build $\Theta(n)$ + $n$ extracts $O(\log n)$ = $\Theta(n\log n)$, space $O(1)$, not stable.

\subsectiontitle{Priority Queue}
\begin{lstlisting}
fn heap_extract_max(a:&mut Vec<i32>)->Option<i32>{
  if a.is_empty(){return None;}
  let ans=a[0]; let last=a.pop().unwrap();
  if !a.is_empty(){ a[0]=last; max_heapify(a,0,a.len()); }
  Some(ans)
}
fn heap_increase_key(a:&mut Vec<i32>, mut i:usize, key:i32){
  if key<a[i]{return;} a[i]=key;
  while i>0 { let p=(i-1)/2; if a[p]>=a[i]{break;} a.swap(i,p); i=p; }
}
fn max_heap_insert(a:&mut Vec<i32>, key:i32){
  a.push(i32::MIN); heap_increase_key(a,a.len()-1,key);
}
\end{lstlisting}
Std min-heap: `BinaryHeap<Reverse<(dist,node)>>`. Ops insert/extract/increase-key $O(\log n)$.

\sectiontitle{Divide \& Conquer Algorithms}
\subsectiontitle{Karatsuba Multiplication}

\[
x=x_hB^m+x_l,\qquad y=y_hB^m+y_l
\]

\[
a=x_hy_h,\quad
b=x_ly_l,\quad
c=(x_h+x_l)(y_h+y_l)
\]

\[
xy=aB^{2m}+(c-a-b)B^m+b
\]

3 recursive multiplications.

3 recursive multiplications: $T(n)=3T(n/2)+\Theta(n)=\Theta(n^{\log_2 3})\approx\Theta(n^{1.585})$.
\subsectiontitle{Maximum Sum Subarray}

Problem:

Input A[0..n) Output max sum contiguous nonempty subarray

D\&C idea:

best subarray is:

\begin{itemize}
  \item entirely left
  \item entirely right
  \item crosses middle
\end{itemize}

\noindent
\begin{minipage}[t]{0.32\linewidth}\vspace{0pt}
\begin{lstlisting}[aboveskip=0pt,belowskip=0pt]
fn max_crossing(a,l,m,r){
  // left suffix [l,m)
  left_sum=-INF;
  sum=0;
  for i in (l..m).rev(){
    sum+=a[i];
    left_sum=max(left_sum,sum);
  }
\end{lstlisting}
\end{minipage}\hfill
\begin{minipage}[t]{0.32\linewidth}\vspace{0pt}
\begin{lstlisting}[aboveskip=0pt,belowskip=0pt]
  // right prefix [m,r)
  right_sum=-INF;
  sum=0;
  for j in m..r {
    sum+=a[j];
    right_sum=max(right_sum,sum);
  }
  left_sum+right_sum
}
\end{lstlisting}
\end{minipage}\hfill
\begin{minipage}[t]{0.32\linewidth}\vspace{0pt}
\begin{lstlisting}[aboveskip=0pt,belowskip=0pt]
fn max_subarray(a,l,r){
  if r-l==1 {return a[l]}
  m=(l+r)/2;
  left=max_subarray(a,l,m);
  right=max_subarray(a,m,r);
  cross=max_crossing(a,l,m,r);
  max(left,max(right,cross))
}
\end{lstlisting}
\end{minipage}

Complexity:

$T(n)=2T(n/2)+\Theta(n)=\Theta(n\log n)$

Linear Kadane:

\begin{lstlisting}[aboveskip=0pt,belowskip=0pt]

fn kadane(a: &Vec<i32>) -> i32 {
    let mut best: i32 = a[0];
    let mut cur: i32 = a[0];

    for i in 1..a.len() {
        cur = a[i].max(cur + a[i]);
        best = best.max(cur);
    }

    best
}

\end{lstlisting}

Kadane invariant:

\[
\text{cur}_i=\max(A_i,\text{cur}_{i-1}+A_i),\qquad \text{best}_i=\max(\text{best}_{i-1},\text{cur}_i)
\]

\subsectiontitle{Matrix Multiplication}

Naive:

\begin{lstlisting}[aboveskip=0pt,belowskip=0pt]
for i in 0..n {
    for j in 0..n {
        for k in 0..n {
            c[i][j] += a[i][k] * b[k][j];
        }
    }
}
\end{lstlisting}

Complexity:

$\Theta(n^3), \text{space} \Theta(n^2)$

 D\&C block formula (8 mults): \texttt{C11=A11B11+A12B21}, \texttt{C12=A11B12+A12B22}, \texttt{C21=A21B11+A22B21}, \texttt{C22=A21B12+A22B22}. Hence $T(n)=8T(n/2)+\Theta(n^2)=\Theta(n^3)$.

 \textbf{Strassen} (7 mults):
\texttt{M1=(A11+A22)(B11+B22)}, \texttt{M2=(A21+A22)B11}, \texttt{M3=A11(B12-B22)}, \texttt{M4=A22(B21-B11)}, \texttt{M5=(A11+A12)B22}, \texttt{M6=(A21-A11)(B11+B12)}, \texttt{M7=(A12-A22)(B21+B22)}.
\texttt{C11=M1+M4-M5+M7}, \texttt{C12=M3+M5}, \texttt{C21=M2+M4}, \texttt{C22=M1-M2+M3+M6}.
\begin{lstlisting}[aboveskip=0pt,belowskip=0pt]
fn strassen(A,B,n){
  if n==1 { return A*B }
  split A,B into A11..A22,B11..B22;
  M1..M7 = strassen(formulas above);
  C11=M1+M4-M5+M7; C12=M3+M5;
  C21=M2+M4;       C22=M1-M2+M3+M6;
  combine(C11,C12,C21,C22)
}
\end{lstlisting}

$T(n)=7T(n/2)+\Theta(n^2)=\Theta(n^{\log_2 7})\approx\Theta(n^{2.807})$.

\sectiontitle{Dynamic Programming}
\subsectiontitle{Rod Cutting}
revenue$[\ell]=$ best revenue length $\ell$, $\text{first}[\ell]=$ first cut.
Rod cutting: revenue[0]=0; revenue[l]=max over 1<=c<=l of price[c]+revenue[l-c].
\begin{lstlisting}[aboveskip=0pt,belowskip=0pt]
fn rod_cut(price:Vec<i32>, n:usize){
  revenue=vec![0;n+1]; first=vec![0;n+1];
  for length in 1..=n {revenue[length]=-INF;
    for cut in 1..=length {
      candidate=price[cut]+revenue[length-cut];
      if candidate>revenue[length] {
        revenue[length]=candidate; first[length]=cut
      }
    }
  }
  (revenue,first)
}
\end{lstlisting}
Reconstruct: output `first[length]`, set `length-=first[length]`. Runtime $\Theta(n^2)$.

\subsectiontitle{Matrix Chain Multiplication}
Matrix $i$: $p_i\times p_{i+1}$. $dp[i][j]=$ min cost for $A_i\cdots A_j$.
MCM recurrence: dp[i][i]=0; dp[i][j]=min over i<=k<j of dp[i][k]+dp[k+1][j]+p[i]*p[k+1]*p[j+1].
\begin{lstlisting}[aboveskip=0pt,belowskip=0pt]
fn matrix_chain(p:Vec<usize>){     // optimal parenthesization
  n=p.len()-1; dp=vec![vec![0;n];n]; split=vec![vec![0;n];n];
  for len in 2..=n {for i in 0..=n-len {
    j=i+len-1; dp[i][j]=INF;
    for k in i..j {
      cost=dp[i][k]+dp[k+1][j]+p[i]*p[k+1]*p[j+1];
      if cost<dp[i][j] {dp[i][j]=cost; split[i][j]=k}
    }
  }}
  (dp,split)
}
\end{lstlisting}
Increasing chain length. $\Theta(n^2)$ states $\times\Theta(n)$ splits $=\Theta(n^3)$.

\subsectiontitle{Optimal BST}
Sorted keys $K_i$ with successful probs $p_i$; optional unsuccessful $q_i$. $e[i][j]=$ min expected cost for subtree $i..j$, $w[i][j]=$ prob mass, $root[i][j]=$ best root.
OBST: e[i][i-1]=0, w[i][i-1]=q[i-1]; e[i][j]=min over root r in [i,j] of e[i][r-1]+e[r+1][j]+w[i][j].
\begin{lstlisting}[aboveskip=0pt,belowskip=0pt]
fn optimal_bst(p,q,n){
  for i in 1..=n+1 {e[i][i-1]=0; w[i][i-1]=q[i-1]}
  for len in 1..=n {for i in 1..=n-len+1 {
    j=i+len-1; e[i][j]=INF; w[i][j]=w[i][j-1]+p[j]+q[j];
    for r in i..=j {
      cost=e[i][r-1]+e[r+1][j]+w[i][j];
      if cost<e[i][j] {e[i][j]=cost; root[i][j]=r}
    }
  }}
  (e,root)
}
\end{lstlisting}
Runtime $\Theta(n^3)$, space $\Theta(n^2)$.

\subsectiontitle{Longest Common Subsequence}
$dp[i][j]=$ LCS length of $A[0..i),B[0..j)$.
LCS: dp[0][*]=dp[*][0]=0. If A[i-1]==B[j-1]: dp[i][j]=1+dp[i-1][j-1]; else max(up,left).
\begin{lstlisting}[aboveskip=0pt,belowskip=0pt]
fn lcs(a:Vec<char>, b:Vec<char>){
  n=a.len(); m=b.len(); dp=vec![vec![0;m+1];n+1];
  for i in 1..=n {for j in 1..=m {
    dp[i][j]=if a[i-1]==b[j-1] {1+dp[i-1][j-1]}
             else {dp[i-1][j].max(dp[i][j-1])}
  }}
  dp[n][m]
}
fn lcs_reconstruct(a:Vec<char>, b:Vec<char>, dp:Vec<Vec<usize>>){
  i=a.len(); j=b.len(); ans=vec![];
  while i>0 && j>0 {
    if a[i-1]==b[j-1] {ans.push(a[i-1]); i-=1; j-=1}
    else if dp[i-1][j]>=dp[i][j-1] {i-=1}
    else {j-=1}
  }
  ans.reverse(); ans
}
\end{lstlisting}
$\Theta(nm)$ time, $\Theta(nm)$ or rolling $\Theta(m)$ space. Equal $\to$ diagonal+1; unequal $\to$ max(up,left).

\subsectiontitle{Knapsack 0/1}
$dp[i][c]=$ best value using first $i$ items, capacity $c$.
0/1 knapsack: dp[i][c]=max(skip dp[i-1][c], take v[i-1]+dp[i-1][c-w[i-1]]) if w[i-1]<=c.
\begin{lstlisting}[aboveskip=0pt,belowskip=0pt]
fn knapsack(w:Vec<usize>, v:Vec<i32>, cap:usize){ // choose item at most once
  n=w.len(); dp=vec![vec![0;cap+1];n+1];
  for i in 1..=n {for c in 0..=cap {
    dp[i][c]=dp[i-1][c];
    if w[i-1]<=c {
      dp[i][c]=dp[i][c].max(v[i-1]+dp[i-1][c-w[i-1]])
    }
  }}
  dp[n][cap]
}
// 1D 0/1: capacity descending. Unbounded: capacity ascending.
\end{lstlisting}
$\Theta(nW)$ pseudo-polynomial.

\subsectiontitle{Coin Change / Change-Making}
Coin change: mincoins[0]=0; mincoins[x]=min over coin<=x of 1+mincoins[x-coin].
\begin{lstlisting}[aboveskip=0pt,belowskip=0pt]
fn coin_change(coins:Vec<usize>, target:usize){ // unlimited coins
  mincoins=vec![INF;target+1]; mincoins[0]=0;   // min number of coins
  for amount in 1..=target {for coin in coins {
    if coin<=amount {
      mincoins[amount]=mincoins[amount].min(1+mincoins[amount-coin])
    }
  }}
  mincoins[target]
}
\end{lstlisting}
Runtime $\Theta(target\cdot \#coins)$.

\subsectiontitle{Longest Increasing Subsequence}
Quadratic:
LIS $O(n^2)$: $dp[i]=1+\max\{dp[j]:j<i, a[j]<a[i]\}$; answer=$\max dp$.
\begin{lstlisting}[aboveskip=0pt,belowskip=0pt]
fn lis_n2(a:Vec<i32>){
  dp=vec![1;a.len()];
  for i in 0..a.len(){for j in 0..i{
    if a[j]<a[i]{dp[i]=dp[i].max(1+dp[j])}
  }}
  max(dp)
}
fn lis_nlogn(a:Vec<i32>){ // tails[len-1] = min tail
  tails=vec![];
  for x in a {
    l=0; r=tails.len();
    while l<r {m=(l+r)/2; if tails[m]<x {l=m+1}else{r=m}}
    if l==tails.len(){tails.push(x)}else{tails[l]=x}
  }
  tails.len()
}
\end{lstlisting}

\subsectiontitle{DAG Longest Path}
Input `adj: WGraph`. Topological order; max-relaxation; negative weights OK because no cycles.
\begin{lstlisting}[aboveskip=0pt,belowskip=0pt]
fn dag_longest(adj:WGraph, source:usize){
  n=adj.len(); dist=vec![-INF;n]; parent=vec![None;n]; dist[source]=0;
  for u in topo(adj) {if dist[u]!=-INF {for (v,w) in adj[u] {
    if dist[u]+w>dist[v] {dist[v]=dist[u]+w; parent[v]=Some(u)}
  }}}
  dist
}
\end{lstlisting}
Runtime $\Theta(V+E)$.
\sectiontitle{Graph Representation / Types}
Use these names everywhere:
\begin{lstlisting}[aboveskip=0pt,belowskip=0pt]
type Graph  = Vec<Vec<usize>>;        // adj[u] = neighbors v
type WGraph = Vec<Vec<(usize,i32)>>;  // adj[u] = (v,weight)
struct Edge { u:usize, v:usize, w:i32 }
type EdgeList = Vec<Edge>;
\end{lstlisting}
\begin{tabular}{@{}lll@{}}
adj list & space $\Theta(V+E)$ & scan neighbors $\Theta(\deg u)$ \\
adj matrix & space $\Theta(V^2)$ & edge lookup $\Theta(1)$, scan $\Theta(V)$ \\
edge list & space $\Theta(E)$ & Bellman-Ford/Kruskal scans all edges
\end{tabular}


\sectiontitle{BFS}
Input `adj: Graph`. Unweighted shortest paths / levels / components / bipartite. Runtime $\Theta(V+E)$.

\noindent
\begin{minipage}[t]{0.49\linewidth}\vspace{0pt}
\begin{lstlisting}[aboveskip=0pt,belowskip=0pt]
fn bfs(adj:Graph, source:usize){  
  n=adj.len(); // shortest path in unweighted graph                    
  dist=vec![INF;n]; // levels, components
  parent=vec![None;n];
  queue=VecDeque::new();
  dist[source]=0;
  queue.push_back(source);

  while let Some(u)=queue.pop_front(){
    for v in adj[u] {
      if dist[v]==INF {
        dist[v]=dist[u]+1;
        parent[v]=Some(u);
        queue.push_back(v);
      }
    }
  }
  (dist,parent)
}
\end{lstlisting}
\end{minipage}\hfill
\begin{minipage}[t]{0.49\linewidth}\vspace{0pt}
\begin{lstlisting}[aboveskip=0pt,belowskip=0pt]
fn bipartite(adj:Graph){          // 2-coloring
  n=adj.len(); color=vec![-1;n];  // false iff odd cycle
  queue=VecDeque::new();
  for s in 0..n {  
    if color[s]!=-1 {continue} 
    color[s]=0; queue.push_back(s);
    while let Some(u)=queue.pop_front(){
      for v in adj[u] {
        if color[v]==-1 {
          color[v]=1-color[u];
          queue.push_back(v);
        }
        else if color[v]==color[u] {return false}
      }
    }
  }
  true
}
\end{lstlisting}
\end{minipage}

\textbf{Invariant:} queue processed in nondecreasing distance. First discovery gives shortest unweighted distance. If a shorter path to $v$ existed, its predecessor would have been processed earlier, so $v$ would already be discovered.

\textbf{Uses:} unweighted shortest path, levels, connected components, bipartite test.

\textbf{Path reconstruction:} follow parent$[t]$ backward until source, then reverse.
\sectiontitle{DFS / Topological Sort / SCC}
Input `adj: Graph`. Colors: WHITE/GRAY/BLACK. Runtime $\Theta(V+E)$.
\begin{lstlisting}[aboveskip=0pt,belowskip=0pt]
fn dfs(adj:Graph){  // topo, cycle detection, SCC
  n=adj.len(); color=vec![WHITE;n]; parent=vec![None;n]; time=0;
  fn visit(u:usize){
    time+=1; discover[u]=time; color[u]=GRAY;
    for v in adj[u] {
      if color[v]==WHITE {parent[v]=Some(u); visit(v)}
    }
    color[u]=BLACK; time+=1; finish[u]=time;
  }
  for u in 0..n {if color[u]==WHITE {visit(u)}}
}
\end{lstlisting}
Directed edges: tree = white target; back = gray ancestor; forward/cross = black. Undirected: only tree/back; cycle test = visited neighbor $\ne$ parent.
DFS times: d[u]<d[v]<f[v]<f[u] => descendant/tree/forward; d[v]<d[u]<f[u]<f[v] => back to ancestor/cycle; disjoint intervals => cross.
Directed cycle detection: edge to GRAY vertex $\Leftrightarrow$ back edge $\Leftrightarrow$ directed cycle. Edge to BLACK does not prove a cycle.
% --- Topological Sort: two minipages, left = topo_dfs, right = topo_kahn ---
\subsectiontitle{Topological Sort}
DAG order: every edge $u\to v$ has $u$ before $v$. DFS = decreasing finish times.
\noindent
\begin{minipage}[t]{0.49\linewidth}\vspace{0pt}
\begin{lstlisting}[aboveskip=0pt,belowskip=0pt]
fn topo_dfs(adj:Graph){
  n=adj.len(); seen=vec![false;n]; order=vec![];
  fn visit(u:usize){
    seen[u]=true;
    for v in adj[u]{if !seen[v]{visit(v)}}
    order.push(u)
  }
  for u in 0..n {if !seen[u]{visit(u)}}
  order.reverse(); order
}
\end{lstlisting}
\end{minipage}\hfill
\begin{minipage}[t]{0.49\linewidth}\vspace{0pt}
\begin{lstlisting}[aboveskip=0pt,belowskip=0pt]
fn topo_kahn(adj:Graph){
  n=adj.len(); indeg=vec![0;n]; 
  queue=VecDeque::new(); order=vec![];
  for u in 0..n {for v in adj[u]{indeg[v]+=1}}
  for u in 0..n {if indeg[u]==0 {queue.push_back(u)}}
  while let Some(u)=queue.pop_front(){
    order.push(u);
    for v in adj[u]{
      indeg[v]-=1;
      if indeg[v]==0{queue.push_back(v)}
    }
  }
  if order.len()==n {Some(order)} else {None}
}
\end{lstlisting}
\end{minipage}
Kahn returns `None` iff not all vertices are output.

% --- Kosaraju: two minipages, split at if seen[source]{continue} ---
\subsectiontitle{Strongly Connected Components}
SCC: $u\leftrightsquigarrow v$. Kosaraju = DFS finish order, reverse graph, DFS in decreasing finish order.
\noindent
\begin{minipage}[t]{0.49\linewidth}\vspace{0pt}
\begin{lstlisting}[aboveskip=0pt,belowskip=0pt]
fn kosaraju(adj:Graph){   // SCC. DFS order -> reverse graph -> DFS
  order=dfs_finish_order(adj); rev=reverse_graph(adj);
  seen=vec![false;n]; comps=[];
  for source in order.rev() {
    if seen[source]{continue}
\end{lstlisting}
\end{minipage}\hfill
\begin{minipage}[t]{0.49\linewidth}\vspace{0pt}
\begin{lstlisting}[aboveskip=0pt,belowskip=0pt]
    comp=[]; stack=[source]; seen[source]=true;
    while let Some(u)=stack.pop(){
      comp.push(u);
      for v in rev[u]{
        if !seen[v]{seen[v]=true; stack.push(v)}
      }
    }
    comps.push(comp)
  }
  comps
}
\end{lstlisting}
\end{minipage}
Runtime $\Theta(V+E)$.

\sectiontitle{Union-Find / Disjoint Sets}

Forest of trees. parent$[x]$ points upward. Root is representative. rank$[x]$ upper-bounds height.

\noindent
\begin{minipage}[t]{0.49\linewidth}\vspace{0pt}
\begin{lstlisting}[aboveskip=0pt,belowskip=0pt]
struct DSU {
  parent: Vec<usize>,
  rank: Vec<usize>,
}
impl DSU{
  fn new(n){DSU{parent:(0..n).collect(),
    rank:vec![0;n]}}

  fn find(&mut self,x){
    if self.parent[x]!=x {
      self.parent[x]=self.find(self.parent[x])
    }
    self.parent[x]
  }
\end{lstlisting}
\end{minipage}\hfill
\begin{minipage}[t]{0.49\linewidth}\vspace{0pt}
\begin{lstlisting}[aboveskip=0pt,belowskip=0pt]
  fn union(&mut self,x,y){
    let (mut rx,mut ry)=(self.find(x),self.find(y));
    if rx==ry {return false}
    if self.rank[rx] < self.rank[ry]
      { swap(rx,ry) }
    self.parent[ry]=rx;
    if self.rank[rx]==self.rank[ry] {
      self.rank[rx]+=1;
    }
    true
  }
}
\end{lstlisting}
\end{minipage}

Complexities:

\begin{itemize}
  \item linked list weighted union: $O(m + n \log n)$
  \item forest + union by rank + path compression: $O(m\alpha(n))$
\end{itemize}

$\alpha(n)\le 5$ in practice.

\noindent
\begin{minipage}[t]{0.49\linewidth}\vspace{0pt}
\begin{lstlisting}[aboveskip=0pt,belowskip=0pt]
//Connected components:
let mut dsu: DSU = DSU::new(n);
for (u, v) in edges {
    dsu.union(u, v);
}
if dsu.find(a) == dsu.find(b) {
    // same component
}
\end{lstlisting}
\end{minipage}\hfill
\begin{minipage}[t]{0.49\linewidth}\vspace{0pt}
\begin{lstlisting}[aboveskip=0pt,belowskip=0pt]
//Cycle detection (undirected):
let mut dsu: DSU = DSU::new(n);
for (u, v) in edges {
    if !dsu.union(u, v) {
        // cycle
    }
}
\end{lstlisting}
\end{minipage}

\sectiontitle{Minimum Spanning Trees}
\subsectiontitle{Chooser}

\begin{tabular}{@{}ll@{}}
sparse graph & Kruskal often easiest \\
dense graph & Prim often easier \\
disconnected graph & forest, no MST
\end{tabular}
tree: connected + acyclic
connected graph:
tree $\Leftrightarrow$ |E|=|V|-1

MST: spanning tree with min total edge weight
\subsectiontitle{Cut / Cycle Properties}
Cut $(S,V-S)$: edge crosses cut iff one endpoint in $S$, one outside. Light edge = minimum-weight crossing edge.

\textbf{Cut property:} if cut respects chosen edge set $A$, then a light crossing edge is safe.

\textbf{Cycle property:} for any cycle, a heaviest edge on that cycle is excluded from some MST. If weights distinct, the heaviest edge is in no MST.

\subsectiontitle{Kruskal}

Data structures:

\begin{lstlisting}[aboveskip=0pt,belowskip=0pt]

struct Edge {
    u,
    v,
    w: i32,
}

let mut edges: Vec<Edge> = ...;
let mut dsu: DSU = DSU::new(n);
let mut mst: Vec<Edge> = Vec::new();
let mut total: i32 = 0;

\end{lstlisting}

Algorithm:

\begin{lstlisting}[aboveskip=0pt,belowskip=0pt]
fn kruskal(n,edges: Vec<Edge>){ // MST
  sort_by_weight(edges); dsu=DSU::new(n); total=0; mst=[];
  for e in edges {              // sort edges + DSU
    if dsu.union(e.u,e.v) {
      total+=e.w; mst.push(e);
      if mst.len()==n-1 {break}
    }
  }
  (total,mst)
}

\end{lstlisting}

Complexity:
sort: $O(E\log E)$
DSU: $O(E\alpha(V))$
total: $O(E\log E)$

 \textbf{Output condition:} if `mst.len()==n-1`, it is an MST; otherwise graph disconnected and result is a forest.

Traps:

\begin{itemize}
  \item If graph disconnected $\to$ MST does not exist; returns forest.
  \item Need sort by weight increasing.
  \item Skip edge if endpoints already connected.
\end{itemize}

\subsectiontitle{Prim}

\begin{lstlisting}[aboveskip=0pt,belowskip=0pt]
let mut adj: Vec<Vec<(usize, i32)>> = vec![Vec::new(); n];
let mut used: Vec<bool> = vec![false; n];
let mut heap: BinaryHeap<Reverse<(i32, usize, usize)>> = BinaryHeap::new();
// (weight, vertex, parent)

\end{lstlisting}

\begin{lstlisting}[aboveskip=0pt,belowskip=0pt]
fn prim(adj,s){                         // MST
  heap.push((0,s)); total=0; cnt=0;     // min crossing edge
  while let Some((w,u))=heap.pop_min(){
    if used[u] {continue}
    used[u]=true; total+=w; cnt+=1;
    for (v,w2) in adj[u] {
      if !used[v] {heap.push((w2,v))}
    }
  }
  if cnt!=n {/* disconnected */}
  total
}
\end{lstlisting}

Complexity:

\begin{itemize}
  \item Binary heap lazy: $O(E\log E)$; for simple graphs $E\le V^2$, so $O(E\log V)$.
\end{itemize}
 \textbf{Output condition:} if `cnt==n`, `total` is MST weight; otherwise graph disconnected.

\subsectiontitle{MST Reduction Stencils}
\textbf{Roads+airports:} add supernode $s$; edge $s$--city cost=airport; old road edges unchanged; MST in augmented graph.
\textbf{Add one new vertex given old MST $T$:} run MST on $T$ plus new incident edges; old non-tree edges still unnecessary by cut/cycle property.
\textbf{Points/routes:} complete graph on points, edge cost=squared distance; run Kruskal/Prim. $E=\Theta(n^2)$, time $O(n^2\log n)$.

\sectiontitle{Shortest Paths}
\subsectiontitle{Chooser}

\begin{tabular}{@{}ll@{}}
unweighted & BFS \\
DAG & topo DP \\
weights $\ge 0$ & Dijkstra \\
negative edges & Bellman-Ford
\end{tabular}
\subsectiontitle{Relaxation}
Relax edge (u,v,w): if dist[u]!=INF and dist[v] > dist[u]+w, set dist[v]=dist[u]+w and parent[v]=u.

\subsectiontitle{Dijkstra}
\begin{lstlisting}[aboveskip=0pt,belowskip=0pt]
fn dijkstra(adj:WGraph, source:usize){ // all weights >= 0
  let n=adj.len();
  let mut dist=vec![INF;n];
  let mut parent=vec![None;n];
  let mut heap=MinHeap::<(i32,usize)>::new();

  dist[source]=0;
  heap.push((0,source));

  while let Some((du,u))=heap.pop(){
    if du!=dist[u] {continue} // stale heap entry

    for (v,w) in adj[u] {
      if dist[v]>du+w {
        dist[v]=du+w;
        parent[v]=Some(u);
        heap.push((dist[v],v));
      }
    }
  }
  dist
}
\end{lstlisting}
Runtime $O((V+E)\log V)$. Requires all weights $\ge0$; negative edges break it.
% --- Bellman-Ford: two minipages, left = init+V-1 passes, right = negative-cycle check+return ---
\subsectiontitle{Bellman-Ford}
Input `edges: EdgeList`. Negative edges OK; detects reachable negative cycles.
\begin{lstlisting}[aboveskip=0pt,belowskip=0pt]
fn bellman_ford(n, edges, source){
  dist=vec![INF;n]; parent=vec![None;n]; dist[source]=0;
  for _ in 0..n-1 { changed=false;
    for e in edges {
      if dist[e.u]!=INF && dist[e.v]>dist[e.u]+e.w {
        dist[e.v]=dist[e.u]+e.w; parent[e.v]=Some(e.u);
        changed=true;
      }
    }
    if !changed {break}
  }
  for e in edges {
    if dist[e.u]!=INF && dist[e.v]>dist[e.u]+e.w {
      return None; // reachable negative cycle
    }
  }
  Some(dist)
}
\end{lstlisting}
Runtime $\Theta(VE)$. Returns `None` iff a reachable negative cycle is detected.
\textbf{Negative cycle not reachable from source:}
Bellman-Ford still returns shortest paths.
Only improvements after pass V reachable from source matter.

\textbf{Weight transform trap:} $w'(e)=a\,w(e)$ with $a>0$ preserves shortest paths. $w'(e)=w(e)+a$ does not preserve shortest paths when candidate paths use different numbers of edges.
% --- DAG shortest path: two minipages, left = init+for u in topo, right = relaxation body+return ---
\subsectiontitle{Shortest Paths in DAG}
\noindent
\begin{minipage}[t]{0.49\linewidth}\vspace{0pt}
\begin{lstlisting}[aboveskip=0pt,belowskip=0pt]
fn dag_shortest(adj:WGraph, source:usize){
  let n=adj.len();                          
  let mut dist=vec![INF;n];                 
  let mut parent=vec![None;n];

  dist[source]=0;

  for u in topo(adj) {
    if dist[u]==INF {continue}
\end{lstlisting}
\end{minipage}\hfill
\begin{minipage}[t]{0.49\linewidth}\vspace{0pt}
\begin{lstlisting}[aboveskip=0pt,belowskip=0pt]
    for (v,w) in adj[u] {    // topological order.
      if dist[v]>dist[u]+w { // one relaxation per edge
        dist[v]=dist[u]+w;   // negative weights ok
        parent[v]=Some(u);
      }
    }
  }
  dist
}
\end{lstlisting}
\end{minipage}
Runtime $\Theta(V+E)$.

\subsectiontitle{Time-Dependent / Ferry Shortest Path}
$dist[u]=$ earliest arrival. Road relax: $dist[v]=\min(dist[v],dist[u]+w)$.
Ferry/event $(u,v,dep,arr)$ usable iff $dist[u]\le dep$, then relax $dist[v]=\min(dist[v],arr)$.
Use PQ by earliest arrival if time is FIFO/nonnegative; otherwise expand state/event graph.

\sectiontitle{Flow Networks}
Directed $G=(V,E)$, source $s$, sink $t$, capacities $c(u,v)\ge0$.
Flow constraints: 0<=f(u,v)<=c(u,v); conservation for u!=s,t: inflow=outflow; |f|=out(s)-in(s).
Augmenting path: $s\to t$ path in residual graph; bottleneck $=\min c_f(e)$ on path.
\textbf{Flow reduction recipe:} selected object = unit flow; capacities = usage limits; source/sink = constraints; value = objective; integral caps => integral solution.
% --- Ford-Fulkerson: two minipages, left = add_edge+bfs_path, right = max_flow ---
\subsectiontitle{Ford-Fulkerson / Edmonds-Karp}
Types: `cap[u][v]` residual capacity matrix; `adj: Graph` stores forward+reverse residual neighbors.
\noindent
\begin{minipage}[t]{0.49\linewidth}\vspace{0pt}
\begin{lstlisting}[aboveskip=0pt,belowskip=0pt]
fn add_edge(u:usize,v:usize,c:i32){
  adj[u].push(v); adj[v].push(u); cap[u][v]+=c
}
fn bfs_path(source:usize, sink:usize){
  parent.fill(None); parent[source]=Some(source);
  queue.push_back((source,INF));
  while let Some((u,flow))=queue.pop_front(){
    for v in adj[u]{
      if parent[v].is_none() && cap[u][v]>0 {
        parent[v]=Some(u);
        pushed=flow.min(cap[u][v]);
        if v==sink {return pushed}
        queue.push_back((v,pushed))
      }
    }
  }
  0
}
\end{lstlisting}
\end{minipage}\hfill
\begin{minipage}[t]{0.49\linewidth}\vspace{0pt}
\begin{lstlisting}[aboveskip=0pt,belowskip=0pt]
fn max_flow(source:usize, sink:usize){ // max flow
  flow=0;               // augment while path exists
  loop {
    pushed=bfs_path(source,sink);
    if pushed==0 {break}
    flow+=pushed;
    v=sink;
    while v!=source {
      u=parent[v].unwrap();
      cap[u][v]-=pushed; cap[v][u]+=pushed;
      v=u
    }
  }
  flow
}
\end{lstlisting}
\end{minipage}
Ford-Fulkerson integer caps: $O(E|f^*|)$. Edmonds-Karp BFS: $O(VE^2)$. Final residual graph gives min cut by residual reachability from $s$.

\subsectiontitle{Max-Flow Min-Cut / Reductions}
Max-flow min-cut: $\max |f|=\min c(S,V\setminus S)$, $s\in S,t\notin S$.
Optimality certificate: flow value $F$ and cut capacity $F$.

Bipartite matching: source$\to L$ cap 1; $L\to R$ cap 1 for original edges; $R\to$sink cap 1. Flow-1 $L\to R$ edges are matched pairs.

Capacitated assignment: $s\to$ person cap 1; person$\to$ acceptable activity cap 1; activity$\to t$ cap capacity. Feasible iff maxflow = \#people. Integral capacities give integral assignment.
\begin{lstlisting}[aboveskip=0pt,belowskip=0pt]
source=0; left0=1; right0=1+n_left; sink=1+n_left+n_right;
for l in 0..n_left {add_edge(source,left0+l,1)}
for (l,r) in bip_edges {add_edge(left0+l,right0+r,1)}
for r in 0..n_right {add_edge(right0+r,sink,1)}
// flow-1 edge left0+l -> right0+r = matched pair
\end{lstlisting}
Edge-disjoint paths: all edges cap 1. Vertex-disjoint: split $v$ into $v_{\text{in}}\to v_{\text{out}}$ cap 1; original $u\to v$ becomes $u_{\text{out}}\to v_{\text{in}}$ cap INF.
\textbf{Grid domino:} checkerboard cells $L/R$, edge for adjacency, max dominoes = max bipartite matching.
\textbf{Path $x\to y$ through $z$:} need two internally vertex-disjoint paths $z\leadsto x$ and $z\leadsto y$; split vertices cap 1, add $x,y\to t$ cap 1, ask maxflow $z\to t$ value 2.
\textbf{Special edge $e=ab$ used:} try orientations $\cdots a-b\cdots$ and $\cdots b-a\cdots$; remove $e$, enforce disjoint prefix/suffix with unit capacities.

\sectiontitle{Randomized / Probabilistic Analysis}
Indicator: $X_A=\mathbf 1[A]$, $\mathbb E[X_A]=\Pr(A)$. Linearity: $E[\sum_i X_i]=\sum_i E[X_i]$.
No independence needed.

Hiring: $P(i\text{ new best})=1/i$, so $E[\#hires]=\sum_i1/i=H_n=\Theta(\log n)$. Hat-check: $P(i\text{ own hat})=1/n$, so $E[\#fixed]=1$.

Randomized quicksort: random pivot avoids input-dependent bad order; expected $\Theta(n\log n)$, worst still $\Theta(n^2)$.

\textbf{Birthday:} $\Pr[\text{all different among }q\text{ into }m]=\prod_{i=0}^{q-1}(1-i/m)\approx e^{-q(q-1)/(2m)}$; for 365 days, $q=23$ gives match prob $>1/2$.
Integral bound for decreasing positive $f$: $\sum_{k=a+1}^{b}f(k)\le\int_a^b f(x)\,dx$.
\sectiontitle{Online Algorithms / Competitive Analysis / Experts}
Online: decide before future known.
Competitive ratio: $ALG(I)\le r\,OPT(I)+c$.
\textbf{Ski rental:} rent 1/day, buy B. Rent until paid B, then buy. If $T\le B$: ALG=T=OPT. If $T>B$: ALG<=B+B=2B=2OPT. Lower bound: adversary stops just before/after purchase; deterministic CR=2.
\subsectiontitle{Caching MC traps}
LRU evicts least-recently requested page. LFU evicts smallest request count; tie usually LRU. Neither has pointwise fewer misses than the other on every sequence. LRU is $k$-competitive for cache size $k$; no deterministic paging algorithm beats $k$ in worst-case competitive ratio.

\subsectiontitle{Experts / Weighted Majority}
$n$ experts, rounds $t=1..T$. Expert $i$ predicts; loss $\ell_{i,t}\in\{0,1\}$.
Weights $w_i=1$ initially. Let $W=\sum_iw_i$.

\begin{lstlisting}[aboveskip=0pt,belowskip=0pt]
for t in 1..=T {
  predict weighted majority of experts;
  observe true outcome;
  for expert i wrong this round {
    w[i] *= beta;        // 0 < beta < 1, often beta=1/2
  }
}
\end{lstlisting}

Weighted prediction: choose label with larger total weight.
Potential: $W_t=\sum_i w_i^{(t)}$.
If ALG makes mistake, wrong-weight $\ge W_t/2$, so
$W_{t+1}\le W_t/2\cdot\beta+W_t/2=W_t(1+\beta)/2$.
If best expert has $m^*$ mistakes, its weight is $\beta^{m^*}$, hence
$\beta^{m^*}\le W_T\le n\big((1+\beta)/2\big)^M$.
For $\beta=1/2$: $M\le 2.41(m^*+\log_2 n)$.
General randomized Hedge: choose expert with prob $w_i/W$; update $w_i\leftarrow w_i e^{-\eta \ell_{i,t}}$; regret $O(\sqrt{T\log n})$.
\subsectiontitle{Prophet Inequality}
Independent nonnegative $X_i$, known distributions. $M=\max_iX_i$, $z=E[M]$, threshold $\tau=z/2$. Accept first $X_i\ge\tau$. Guarantee $E[ALG]\ge z/2$. Proof stencil: $q=P[\forall i:X_i<\tau]$, $b_i=E[(X_i-\tau)^+]$. Value from threshold part $\ge\tau(1-q)$; overshoot part $\ge q\sum_i b_i$; combine to get $\ge z/2$.

Secretary: random order; observe first $n/e$ candidates, remember best rank seen, then hire first later candidate better than all observed.
\[
P(\text{success})\approx 1/e
\]

\sectiontitle{WORKED EXAM PROBLEMS}
{\color{rustMacro}\itshape Representative problems from exercise sets / exams. Question $\to$ exact answer skeleton.}\par\vspace{0.006cm}

\workedq{Q1 (Ex1)}{Prove Linear-Search correct with a loop invariant.}{
\inv at start of iteration $i$, $A[j]\ne v$ for all $j<i$. \init $i=1$: empty range. \maint if invariant true and $A[i]\ne v$, then $A[j]\ne v$ for all $j<i+1$. \term either returns $i$ with $A[i]=v$, or $i=n+1$ so $A[j]\ne v$ for all $j$ and returns NIL.
}

\workedq{Q2 (Ex1)}{Show $(n+a)^b=\Theta(n^b)$ for constants $a,b>0$.}{
For $n\ge2|a|$: $\frac12n\le n+a\le2n$, hence $(\frac12)^bn^b\le(n+a)^b\le2^bn^b$. Both sides are constants times $n^b$; therefore $(n+a)^b=\Theta(n^b)$.
}

\workedq{Q3 (Ex2)}{Solve $T(n)=8T(n/2)+\Theta(n^2)$ by recursion tree.}{
Level $i$: $8^i$ nodes, each cost $c(n/2^i)^2$, total $cn^2 2^i$. Sum $\sum_{i=0}^{\lg n}cn^2 2^i=\Theta(n^2\cdot n)=\Theta(n^3)$. Leaves: $n^{\lg8}=n^3$. Hence $T(n)=\Theta(n^3)$.
}

\workedq{Q4 (Ex2)}{Recursive insertion sort recurrence.}{
$T(n)=T(n-1)+\Theta(n)$. Unroll: $T(n)=c\sum_{i=1}^n i=c\,n(n+1)/2=\Theta(n^2)$.
}

\workedq{Q5 (Ex6)}{$n$ restaurants on a line, positions $M_i$, profit $p_i$, any two at distance $\ge k$. Max profit.}{
Sort by position. State $dp[i]=$ best profit using restaurants $i..n$. Let $next(i)=$ first $j>i$ with $M_j\ge M_i+k$. Base $dp[n+1]=0$. Recurrence $dp[i]=\max(dp[i+1],p_i+dp[next(i)])$. Order $i=n..1$. \ans $dp[1]$. Runtime $O(n\log n)$ with binary search, $O(n^2)$ if scanning. Correctness: optimum either skips $i$ or takes $i$; after taking $i$, remaining choices start at $next(i)$.
}

\workedq{Q5b (Ex6)}{Restaurant leftmost-open variant.}{
Sort positions. $c[i]=$ best total profit if $i$ is leftmost opened. Then $c[i]=p_i+\max\{c[j]:j>i\land M_j\ge M_i+k\}$. \ans $\max_i c[i]$. Bottom-up $O(n^2)$; with binary search/suffix maxima $O(n\log n)$. Fixing the leftmost open restaurant leaves an independent feasible suffix.
}

\workedq{Q6 (Ex5)}{$O(\log n)$ test whether sorted $A[1..n]$ contains $k$.}{
Binary search. Invariant: if $k$ occurs, then $k\in A[l..r)$. Compare $k$ with $A[mid]$ and keep only the half that can contain $k$. Sortedness preserves invariant. Termination: range empty $\Rightarrow$ absent, or equality found. Runtime $T(n)=T(n/2)+O(1)=O(\log n)$.
}

\workedq{Q7 (Ex9)}{Maximum bipartite matching as a flow.}{
Build network: add $s\to L$, original edges $L\to R$, $R\to t$, all capacities 1. Integral max-flow gives integral solution. Capacity 1 forces each left/right vertex used at most once. Unit-flow edge $l\to r$ is a matched pair; every matching gives same-size flow. Thus max flow = max matching.
}


\workedq{Q9 (Ex12)}{What does Lomuto Partition return on an all-equal array, and how to fix it?}{
With condition $\le pivot$, every element is moved left, so $q=r$. Split sizes $n-1$ and 0, hence $T(n)=T(n-1)+\Theta(n)=\Theta(n^2)$. Fix: 3-way partition into $<,=,>$ pivot, or split equal keys evenly / return middle of equal block.
}

\workedq{Q10 (L19 exam)}{Ski resort: lifts/slopes have signed altitude changes. Verify every closed tour returns to start with net change 0.}{
Build digraph: stations = vertices, lift/slope = directed edge weighted by altitude change. Bad tour $\Leftrightarrow$ directed cycle with nonzero total weight. Run Bellman-Ford from super-source on weights $w$ and on $-w$: a positive cycle in $w$ is negative in $-w$; a negative cycle in $w$ is detected in $w$. Relaxation on pass $V$ in either run $\Rightarrow$ nonzero cycle exists. Runtime $O(VE)$.
}

\workedq{Q11 (L19)}{(a) Min crossing edge of any cut is in the MST. (b) Max edge of any cycle is not.}{
Assume distinct weights. (a) Cut property: if lightest crossing edge $e$ absent, add $e$ to MST; cycle contains another crossing edge $f$ with $w(f)>w(e)$; swap $f\mapsto e$ gives cheaper tree, contradiction. (b) If heaviest cycle edge $e$ in MST, remove it; the cycle has lighter edge $f$ crossing the induced cut; swap $e\mapsto f$ gives cheaper tree, contradiction.
}

\workedq{Q12}{Travel passes DP.}{
Travel days $d_1<\cdots<d_n$, pass lengths $L\in\{1,7,30\}$, costs $c_L$. State $dp[i]=$ min cost to cover $d_i..d_n$. Base $dp[n+1]=0$. Let $next(i,L)=$ first $j$ with $d_j\ge d_i+L$. Recurrence $dp[i]=\min_{L\in\{1,7,30\}}(c_L+dp[next(i,L)])$. Order $i=n..1$. \ans $dp[1]$. Runtime $O(n\log n)$ with lower\_bound. Trap: pass covers days $<d_i+L$, not $\le d_i+L$.
}

\workedq{Q13}{Weighted interval scheduling.}{
Sort intervals by finish time. $p(i)=$ last interval finishing $\le start_i$. State $dp[i]=$ best weight using first $i$ intervals. Base $dp[0]=0$. Recurrence $dp[i]=\max(dp[i-1],w_i+dp[p(i)])$. \ans $dp[n]$. Runtime $O(n\log n)$ with binary search for $p(i)$. Correctness: optimum either excludes interval $i$ or includes it, then previous chosen interval must be $\le p(i)$.
}

\workedq{Q14}{Interval scheduling greedy.}{
Sort by increasing finish time. Take interval if $start\ge last\_finish$. Proof: let $e$ be earliest finishing interval. For any optimum $O$, replace the first interval of $O$ by $e$; finish time does not increase, so all later intervals remain feasible. Therefore some optimum contains $e$. Induct on remaining intervals.
}

\workedq{Q15}{DP design answer skeleton.}{
State $dp[s]=$ optimum value for subproblem $s$. Base: smallest/empty cases. Recurrence $dp[s]=\operatorname{opt}_{c\in valid(s)}F(s,c,dp[next(s,c)])$. Order: dependencies first. \ans explicitly write returned state. Runtime = \#states $\times$ transition cost. Correctness: induction in fill order; recurrence enumerates all valid choices; dependencies optimal by IH.
}

\workedq{Q16a}{Parametric search / split array into $k$ contiguous parts minimizing max load.}{
Bounds $L=\max_i a_i$, $R=\sum_i a_i$. feasible$(M)$: scan left-to-right, start new part when current sum $+a_i>M$. Count parts. If count $\le k$, $M$ feasible. Binary search smallest feasible $M$. Correctness: greedy uses minimum possible number of parts for fixed $M$; feasibility monotone. Runtime $O(n\log(\sum a_i))$.
}

\workedq{Q16b}{Forbidden mutual reachability pairs in directed graph.}{
Compute SCCs in $O(V+E)$. For each forbidden pair $(u,v)$, reject iff $comp[u]=comp[v]$. If same SCC, then $u\leadsto v$ and $v\leadsto u$. If mutually reachable, they must be in same SCC. Thus all constraints checked exactly.
}

\workedq{Q16c}{Capacitated activity assignment.}{
Network: $s\to$ each student cap 1; student$\to$ acceptable activity cap 1; activity$\to t$ cap number of spots. Run maxflow. 
\ans feasible iff flow value = \#students. Correctness: integral unit leaving each student chooses one acceptable activity; activity capacity enforces spots; any assignment gives same flow.
}

\workedq{Q16d}{Binary tree diameter / longest path.}{
Postorder. Empty height $=-1$. For node $u$: $h(u)=1+\max(h_L,h_R)$ and $diam(u)=\max(diam_L,diam_R,h_L+h_R+2)$. 
\ans $diam(root)$. Correctness: every longest path either stays inside one child subtree or crosses through $u$ using deepest node on each side.
}

\workedq{Q16e}{Connected graph with $n$ vertices and $n$ edges: count spanning trees.}{
Such graph has exactly one cycle. Find the cycle by DFS back edge or by repeatedly deleting degree-1 leaves. If cycle length is $c$, answer is $c$. Reason: a spanning tree has $n-1$ edges, so delete one edge. Deleting a cycle edge keeps connected and removes the only cycle; deleting a non-cycle edge disconnects.
}

\workedq{Q16f}{Coin-row 2-player minimax DP.}{
$dp[i][j]=$ max value current player can guarantee from coins $i..j$. Base $dp[i][i]=a_i$. Recurrence $dp[i][j]=\max(a_i+\min(dp[i+2][j],dp[i+1][j-1]),\ a_j+\min(dp[i+1][j-1],dp[i][j-2]))$. Opponent chooses the branch minimizing our future gain. Fill by interval length.
}

\workedq{Q17}{Competitive-ratio template.}{
Define OPT$(I)$. Compute ALG$(I)$. Split into exhaustive cases. In each case prove ALG$(I)\le r\,OPT(I)+c$. Take worst case. Conclusion: algorithm is $r$-competitive. Lower bound: adversary chooses future after seeing deterministic algorithm action.
}

\workedq{Q18}{Secretary problem.}{
Random order. Reject first $\lfloor n/e\rfloor$ candidates, remember best seen, then hire first later candidate better than all previous. Success probability tends to $1/e$. Reason: best candidate must appear after sample, and best candidate before it must lie inside sample.
}

\workedq{Q19}{Online bipartite matching.}{
Left side known, right vertices arrive online. Greedy: match arriving vertex to any free neighbor. Each greedy edge can be charged by at most two OPT edges, so ALG $\ge\frac12$OPT. Ranking: random rank/permutation on left vertices; match arriving vertex to lowest-ranked free neighbor. Expected $(1-1/e)$OPT.
}
\workedq{Q20}{Prophet inequality threshold algorithm.}{
Known independent values $X_i$. Let $M=\max_iX_i$, $z=E[M]$, threshold $\tau=z/2$. Accept first value at least $\tau$. Guarantee $E[ALG]\ge z/2$. Write proof with $q=P[\text{no value reaches }\tau]$ and overshoots $(X_i-\tau)^+$; threshold contribution plus overshoot contribution gives half of prophet value.
}

\workedq{Q21}{Inorder+postorder tree reconstruction.}{
root=last(postorder). Split inorder at root into left/right. Sizes determine postorder left/right blocks. Recurse. Correctness by induction on subtree size. Runtime $O(n)$ with key$\to$inorder-index map.
}

\workedq{Q22}{Longest path passing through root.}{
Compute subtree heights. Answer $h(left)+h(right)+2$ with empty height $-1$. Correct because path through root must be deepest node in left side + root + deepest node in right side.
}

\workedq{Q23}{Domino/grid matching.}{
Checkerboard color cells. Put black cells in $L$, white cells in $R$, edge iff adjacent. A domino is one matched edge; no cell reused because matching. Max domino placement = max bipartite matching.
}

\end{multicols*}

\end{document}
